\documentclass[a4paper]{article}

\usepackage[utf8]{inputenc}
\usepackage[T5]{fontenc}
\usepackage[vietnamese]{babel}
\usepackage[fontsize=13pt]{scrextend}
\usepackage{mathptmx}
\usepackage{amsmath}
\usepackage{geometry}
\usepackage{setspace}
\usepackage{indentfirst}
\usepackage{graphicx}
\usepackage{xcolor}
\usepackage{booktabs}
\usepackage{longtable}
\usepackage{tabularx}
\usepackage{array}
\usepackage{enumitem}
\usepackage{listings}
\usepackage{caption}
\usepackage{fancyhdr}
\usepackage{hyperref}
\usepackage{microtype}
\usepackage{tikz}
\usetikzlibrary{arrows.meta,positioning,fit}

\geometry{
  left=3cm,
  right=2cm,
  top=2cm,
  bottom=2cm
}
\setstretch{1.3}
\setlength{\parindent}{1cm}
\setlength{\parskip}{3pt}
\setlength{\emergencystretch}{3em}
\setlist{nosep}
\renewcommand{\arraystretch}{1.25}

\definecolor{codebg}{RGB}{247,248,250}
\definecolor{codeframe}{RGB}{210,214,220}
\definecolor{codekeyword}{RGB}{34,87,153}
\definecolor{codecomment}{RGB}{60,130,80}
\definecolor{codestring}{RGB}{170,70,50}
\definecolor{linkblue}{RGB}{20,80,160}

\hypersetup{
  colorlinks=true,
  linkcolor=linkblue,
  urlcolor=linkblue,
  citecolor=linkblue,
  pdftitle={Báo cáo mã nguồn MSMS Project},
  pdfauthor={Nhóm phát triển MSMS}
}

\lstdefinestyle{cstyle}{
  language=C,
  backgroundcolor=\color{codebg},
  frame=single,
  rulecolor=\color{codeframe},
  basicstyle=\ttfamily\small,
  keywordstyle=\color{codekeyword}\bfseries,
  commentstyle=\color{codecomment}\itshape,
  stringstyle=\color{codestring},
  numbers=left,
  numberstyle=\scriptsize\color{gray},
  numbersep=8pt,
  breaklines=true,
  breakatwhitespace=false,
  showstringspaces=false,
  tabsize=2,
  columns=fullflexible,
  keepspaces=true,
  captionpos=b
}

\lstdefinestyle{textstyle}{
  backgroundcolor=\color{codebg},
  frame=single,
  rulecolor=\color{codeframe},
  basicstyle=\ttfamily\small,
  breaklines=true,
  columns=fullflexible,
  showstringspaces=false
}

\pagestyle{fancy}
\fancyhf{}
\fancyhead[L]{Báo cáo mã nguồn MSMS\_Project}
\fancyhead[R]{ESP-IDF}
\fancyfoot[C]{\thepage}
\setlength{\headheight}{16pt}

\newcommand{\projectpath}[1]{\texttt{\detokenize{#1}}}
\newcommand{\codeword}[1]{\texttt{\detokenize{#1}}}

\begin{document}
\pagenumbering{Alph}

\begin{titlepage}
  \centering
  \vspace*{1.2cm}
  {\large BÁO CÁO PHÂN TÍCH MÃ NGUỒN\par}
  \vspace{1.5cm}
  {\Huge\bfseries MSMS\_Project\par}
  \vspace{0.8cm}
  {\Large Firmware nút cảm biến trong mạng Wi-Fi Mesh\par}
  \vspace{2cm}

  \begin{tabular}{rl}
    Nền tảng: & ESP32 / ESP-IDF / FreeRTOS \\
    Ngôn ngữ: & C, CMake \\
    Kiến trúc mạng: & ESP-Mesh-Lite, Wi-Fi, UDP, UART \\
    Phiên bản tài liệu: & 1.0 \\
  \end{tabular}

  \vfill
  {\large Tài liệu được xây dựng từ trạng thái mã nguồn hiện tại\par}
  \vspace{0.3cm}
  {\large Tháng 07 năm 2026\par}
\end{titlepage}

\pagenumbering{roman}
\tableofcontents
\clearpage
\listoftables
\clearpage
\pagenumbering{arabic}

\section{Giới thiệu}

\subsection{Mục đích tài liệu}

Tài liệu này mô tả cấu trúc, luồng thực thi và các giao diện quan trọng của
\codeword{MSMS_Project}. Mục tiêu là giúp người phát triển:

\begin{itemize}
  \item hiểu quá trình khởi động firmware và quan hệ giữa các thành phần;
  \item xác định nơi thêm cảm biến, giao diện hoặc phương thức truyền dữ liệu;
  \item hiểu định dạng telemetry từ node tới gateway;
  \item có cơ sở kiểm thử, gỡ lỗi và bảo trì hệ thống.
\end{itemize}

Phạm vi tài liệu chỉ bao gồm firmware trong thư mục
\projectpath{MSMS_Project}. Dashboard, gateway bên ngoài và firmware của các
vi điều khiển khác chỉ được nhắc đến tại điểm giao tiếp với dự án này.

\subsection{Vai trò của firmware}

\codeword{MSMS_Project} điều khiển một nút cảm biến dùng ESP32. Firmware đọc
cảm biến tại ba cổng, hiển thị trạng thái trên OLED, quản lý pin, thời gian và
lưu trữ, sau đó đóng gói dữ liệu thành JSON để gửi qua mạng mesh. Một thiết bị
có thể chuyển sang vai trò root; khi đó dữ liệu UDP nhận từ các node được đưa
vào hàng đợi và chuyển tiếp qua UART tới gateway.

\section{Nền tảng và tổ chức dự án}

\subsection{Công nghệ sử dụng}

\begin{table}[ht]
  \centering
  \caption{Các công nghệ chính}
  \begin{tabularx}{\textwidth}{>{\bfseries}p{3.5cm}X}
    \toprule
    Thành phần & Vai trò \\
    \midrule
    ESP-IDF & SDK, driver ngoại vi, Wi-Fi, NVS, SPIFFS và hệ thống build. \\
    FreeRTOS & Task, queue, mutex, timer và điều phối thời gian thực. \\
    ESP-Mesh-Lite & Hình thành mạng nhiều tầng và chuyển tiếp IP giữa các node. \\
    CMake & Khai báo project và các component nội bộ. \\
    cJSON & Phân tích và tái cấu trúc gói JSON tại root/gateway. \\
    NVS, SPIFFS, SD card & Lưu cấu hình, tài nguyên web và dữ liệu cục bộ. \\
    \bottomrule
  \end{tabularx}
\end{table}

\subsection{Cấu trúc thư mục}

\begin{longtable}{p{5cm}p{9cm}}
  \caption{Cấu trúc mã nguồn chính}\\
  \toprule
  \textbf{Đường dẫn} & \textbf{Chức năng} \\
  \midrule
  \endfirsthead
  \toprule
  \textbf{Đường dẫn} & \textbf{Chức năng} \\
  \midrule
  \endhead
  \projectpath{main/} & Điểm vào \codeword{app_main()}, khởi tạo toàn hệ thống. \\
  \projectpath{component/core/} & Dữ liệu trung tâm, pin, thời gian, lưu trữ,
  hiệu năng, xử lý mesh và UART gateway. \\
  \projectpath{component/sensors/} & Kiểu cảm biến, registry và lớp quản lý
  cảm biến. \\
  \projectpath{component/drivers/} & Driver phần cứng: BME280, DHT, AHT10,
  PMS7003, MH-Z14A, DS3231, SSD1306, PCF8574 và các driver hỗ trợ. \\
  \projectpath{component/ui/} & Nút nhấn, màn hình và hệ thống menu. \\
  \projectpath{component/network/} & InternetManager, Wi-Fi, MeshManager,
  WebSocket và trang cấu hình Wi-Fi. \\
  \projectpath{component/utils/} & Mã lỗi và các tiện ích bit. \\
  \projectpath{openwiki/} & Tài liệu kiến trúc và hướng dẫn phát triển. \\
  \projectpath{partitions.csv} & Bảng phân vùng flash cho NVS, OTA và SPIFFS. \\
  \bottomrule
\end{longtable}

File CMake cấp dự án bổ sung sáu nhóm component vào
\codeword{EXTRA_COMPONENT_DIRS}. Ảnh SPIFFS được tạo từ
\projectpath{component/network/WebConfigWifi} và được flash cùng firmware.

\section{Kiến trúc phần mềm}

\subsection{Mô hình phân lớp}

Hệ thống được tổ chức theo hướng component. Driver chỉ xử lý phần cứng;
lớp sensor chuẩn hóa giao diện cảm biến; core điều phối dữ liệu dùng chung;
network chịu trách nhiệm kết nối; UI đọc và cập nhật trạng thái thông qua
\codeword{DataManager_t}.

\begin{figure}[ht]
  \centering
  \begin{tikzpicture}[
    layer/.style={draw=linkblue,rounded corners=3pt,minimum width=13.5cm,
      minimum height=1.05cm,align=center,font=\small},
    flow/.style={-{Latex[length=2.5mm]},thick,draw=linkblue}
  ]
    \node[layer,fill=blue!8] (app)
      {\textbf{Lớp ứng dụng}\\
       \codeword{app_main}, MenuSystem, hành vi vận hành thiết bị};
    \node[layer,fill=cyan!9,below=4mm of app] (core)
      {\textbf{Lớp lõi}\\
       DataManager, Battery, Time, Storage, SystemPerformance,
       ProcessingDataMesh};
    \node[layer,fill=green!9,below=4mm of core] (service)
      {\textbf{Lớp dịch vụ}\\
       SensorManager và Registry \quad InternetManager, Mesh, Wi-Fi
       \quad Button/Screen};
    \node[layer,fill=orange!10,below=4mm of service] (driver)
      {\textbf{Lớp driver}\\
       I2C, UART, GPIO, OLED, RTC, SD card và driver từng cảm biến};
    \node[layer,fill=gray!12,below=4mm of driver] (platform)
      {\textbf{Nền tảng}\\ ESP-IDF, FreeRTOS và phần cứng ESP32};

    \draw[flow] (app) -- (core);
    \draw[flow] (core) -- (service);
    \draw[flow] (service) -- (driver);
    \draw[flow] (driver) -- (platform);
  \end{tikzpicture}
  \caption{Kiến trúc phân lớp của MSMS\_Project}
\end{figure}

\paragraph{Ý nghĩa các lớp.}
Mũi tên đi từ trên xuống thể hiện lớp phía trên sử dụng dịch vụ của lớp phía
dưới. Lớp ứng dụng quyết định trình tự khởi động và hành vi thiết bị nhưng
không truy cập trực tiếp thanh ghi phần cứng. Lớp lõi giữ trạng thái và điều
phối tác vụ. Lớp dịch vụ cung cấp các chức năng cảm biến, mạng và giao diện.
Lớp driver chuyển yêu cầu chung thành thao tác I2C, UART, GPIO hoặc SPI. Cuối
cùng, ESP-IDF và FreeRTOS cung cấp API hệ thống, scheduler và driver nền tảng.

\subsection{Sơ đồ khối toàn hệ thống}

Sơ đồ khối dưới đây mô tả ranh giới giữa phần cứng, firmware trên ESP32 và các
thành phần bên ngoài. Khối màu xanh nằm trong firmware
\codeword{MSMS_Project}; các khối màu vàng là thiết bị vật lý kết nối trực
tiếp; các khối màu tím thuộc tuyến gateway và hệ thống giám sát.

\begin{figure}[ht]
  \centering
  \begin{tikzpicture}[
    hw/.style={draw=orange!80!black,fill=orange!12,rounded corners=3pt,
      minimum width=2.6cm,minimum height=0.9cm,align=center,font=\scriptsize},
    fw/.style={draw=linkblue,fill=blue!8,rounded corners=3pt,
      minimum width=3.1cm,minimum height=0.95cm,align=center,font=\scriptsize},
    ext/.style={draw=violet!75!black,fill=violet!9,rounded corners=3pt,
      minimum width=2.7cm,minimum height=0.95cm,align=center,font=\scriptsize},
    arrow/.style={-{Latex[length=2.2mm]},thick,draw=gray!75},
    bus/.style={<->,thick,draw=gray!75},
    group/.style={draw=linkblue,dashed,rounded corners=4pt,inner sep=4mm}
  ]
    \node[hw] (sensors) {Cảm biến môi trường\\I2C/UART/Pulse};
    \node[hw,below=5mm of sensors] (rtc) {RTC DS3231};
    \node[hw,below=5mm of rtc] (battery) {Pin và mạch đo ADC};
    \node[hw,below=5mm of battery] (buttons) {Nút nhấn};
    \node[hw,below=5mm of buttons] (oled) {OLED SSD1306};
    \node[hw,below=5mm of oled] (sd) {SD card};

    \node[fw,right=13mm of sensors] (sensorSvc)
      {Sensor Registry\\Sensor Manager};
    \node[fw,below=5mm of sensorSvc] (timeSvc)
      {Time Manager\\Battery Manager};
    \node[fw,below=5mm of timeSvc] (dm)
      {\textbf{DataManager}\\trạng thái dùng chung};
    \node[fw,below=5mm of dm] (ui)
      {Menu System\\Screen/Button};
    \node[fw,below=5mm of ui] (storage)
      {Storage Manager\\NVS/SPIFFS/SD};

    \node[fw,right=13mm of sensorSvc] (mesh)
      {MeshManager\\InternetManager};
    \node[fw,below=5mm of mesh] (process)
      {ProcessingDataMesh\\packet loss theo MAC};
    \node[fw,below=5mm of process] (uart)
      {UartToGateWay\\RX lệnh / TX JSON};
    \node[fw,below=5mm of uart] (perf)
      {SystemPerformance\\ErrorCodes};

    \node[ext,right=14mm of mesh] (otherNode)
      {Các mesh node\\khác};
    \node[ext,below=8mm of otherNode] (gateway)
      {Gateway ngoài};
    \node[ext,below=8mm of gateway] (backend)
      {Backend / Dashboard};

    \draw[bus] (sensors) -- node[above,font=\tiny]{driver} (sensorSvc);
    \draw[bus] (rtc) -- (timeSvc);
    \draw[arrow] (battery) -- (timeSvc);
    \draw[arrow] (buttons) -- (ui);
    \draw[arrow] (ui) -- (oled);
    \draw[bus] (storage) -- (sd);

    \draw[arrow] (sensorSvc) -- (dm);
    \draw[arrow] (timeSvc) -- (dm);
    \draw[bus] (ui) -- (dm);
    \draw[bus] (storage) -- (dm);
    \draw[arrow] (perf) -- (dm);
    \draw[bus] (dm) -- (mesh);
    \draw[arrow] (mesh) -- (process);
    \draw[arrow] (process) -- (uart);
    \draw[bus] (mesh) -- node[above,font=\tiny]{Wi-Fi Mesh/UDP} (otherNode);
    \draw[arrow] (uart) -- node[above,font=\tiny]{UART} (gateway);
    \draw[arrow] (gateway) -- node[right,font=\tiny]{WebSocket/API} (backend);

    \node[group,fit=(sensorSvc)(storage)(mesh)(perf),
      label={[font=\small\bfseries]above:Firmware MSMS\_Project}] {};
  \end{tikzpicture}
  \caption{Sơ đồ khối phần cứng và phần mềm của hệ thống}
\end{figure}

Trong chế độ node, đường dữ liệu chính đi từ cảm biến qua SensorManager,
DataManager rồi tới MeshManager để phát UDP. Trong chế độ root, MeshManager
nhận UDP từ các node khác, ProcessingDataMesh chuẩn hóa dữ liệu và task UART
chuyển từng dòng JSON sang gateway. Các khối UI, lưu trữ, pin và hiệu năng hoạt
động song song, cùng cập nhật hoặc đọc trạng thái trong DataManager.

\subsection{DataManager là trạng thái trung tâm}

\codeword{DataManager_t} là đối tượng dùng chung giữa các component. Cấu trúc
này chứa:

\begin{itemize}
  \item dữ liệu cảm biến và cảm biến được gán cho từng cổng;
  \item trạng thái nút nhấn, màn hình, Wi-Fi và mesh;
  \item callback, thông tin phiên bản firmware và mảng mã lỗi;
  \item handle của các task do hệ thống quản lý;
  \item ngữ cảnh mesh gồm queue, mutex, role, socket UDP và địa chỉ đích.
\end{itemize}

Thiết kế này giúp các module trao đổi trạng thái thuận tiện, nhưng cũng tạo
ra rủi ro truy cập đồng thời. Khi bổ sung trường được nhiều task ghi, cần dùng
mutex, queue hoặc vùng tới hạn thay vì đọc/ghi trực tiếp không đồng bộ.

\section{Luồng khởi động}

\subsection{Điểm vào app\_main}

Trình tự khởi động trong \projectpath{main/main.c} như sau:

\begin{enumerate}
  \item khởi tạo NVS; nếu phân vùng hết trang hoặc sai phiên bản thì xóa và
  khởi tạo lại;
  \item mount SPIFFS tại \projectpath{/spiffs};
  \item khởi tạo nút nhấn và bus I2C dùng chung;
  \item tạo driver OLED SSD1306 và khởi tạo ScreenManager;
  \item đặt toàn bộ ba cổng về \codeword{SENSOR_NONE}, sau đó khởi tạo menu;
  \item khởi tạo UART nối gateway;
  \item khởi tạo BatteryManager và task đo pin;
  \item khởi tạo SystemPerformance;
  \item khởi tạo InternetManager ở chế độ mesh, vai trò ban đầu là node;
  \item trên ESP32, khởi tạo RTC DS3231 và StorageManager;
  \item tạo hai task điều hướng và render menu.
\end{enumerate}

\begin{lstlisting}[style=cstyle,caption={Phần cốt lõi của quá trình khởi động}]
void app_main(void)
{
    nvs_flash_init();
    init_spiffs();
    ButtonManagerInit(&DataManager);
    i2cInitDevCommon();

    ScreenManagerInit(&DataManager, oled_dev);
    MenuSystemInit(&DataManager);
    UartToGateWay_Init(&DataManager);

    BatteryManager_Init(&DataManager, ADC_UNIT_1);
    BatteryManager_StartTask(&DataManager);
    SystemPerformance_Init(&DataManager);

    InternetManager_Init(&DataManager, INTERNET_MODE_MESH);
}
\end{lstlisting}

Mặc định firmware không khởi động ở chế độ Wi-Fi độc lập. Muốn thay đổi cần
gọi \codeword{InternetManager_SwitchMode()} hoặc sửa tham số trong
\codeword{app_main()}.

\section{Hệ thống cảm biến}

\subsection{Giao diện driver thống nhất}

Mỗi cảm biến được mô tả bởi \codeword{sensor_driver_t}, gồm tên, mô tả,
đơn vị đo, số lượng giá trị, loại giao tiếp và ba callback
\codeword{init/read/deinit}. Nhờ đó SensorManager và MeshManager không cần biết
chi tiết thanh ghi hay giao thức của từng thiết bị.

\begin{lstlisting}[style=cstyle,caption={Giao diện chung của driver cảm biến}]
typedef struct {
    const char *name;
    char description[20];
    char unit[20];
    bool is_init;
    uint8_t unit_count;
    SensorInterface interface;
    system_err_t (*init)(SensorPort_t *port);
    system_err_t (*read)(SensorPort_t *port, SensorData_t *data);
    system_err_t (*deinit)(SensorPort_t *port);
} sensor_driver_t;
\end{lstlisting}

\subsection{Các cảm biến đã đăng ký}

\begin{table}[ht]
  \centering
  \caption{Danh sách cảm biến trong SensorRegistry}
  \begin{tabularx}{\textwidth}{c l l X}
    \toprule
    ID & Cảm biến & Giao tiếp & Dữ liệu chính \\
    \midrule
    0 & BME280 & I2C & Nhiệt độ, áp suất, độ ẩm \\
    1 & MH-Z14A & UART & Nồng độ CO$_2$, nhiệt độ \\
    2 & PMS7003 & UART & PM1.0, PM2.5, PM10 \\
    3 & DHT22 & Pulse & Nhiệt độ, độ ẩm \\
    4 & AHT10 & I2C & Nhiệt độ, độ ẩm \\
    5 & DHT11 & Pulse & Nhiệt độ, độ ẩm \\
    6 & HTU21D & I2C & Nhiệt độ, độ ẩm \\
    \bottomrule
  \end{tabularx}
\end{table}

\subsection{Quy trình thêm cảm biến}

\begin{enumerate}
  \item tạo hoặc hoàn thiện driver trong \projectpath{component/drivers/};
  \item thêm giá trị vào \codeword{SensorType} mà không thay đổi ID đã phát hành;
  \item khai báo một \codeword{sensor_driver_t} với đủ callback;
  \item thêm driver vào bảng trong \projectpath{SensorRegistry.c};
  \item cập nhật CMake và menu lựa chọn cảm biến nếu cần;
  \item kiểm thử khởi tạo, đọc, mất kết nối và khởi tạo lại;
  \item cập nhật phía gateway để hiểu ID và thứ tự các giá trị mới.
\end{enumerate}

ID cảm biến là một phần của giao thức truyền dữ liệu. Việc chèn một enum vào
giữa danh sách có thể khiến dashboard diễn giải sai dữ liệu đã triển khai.

\section{Mạng và truyền dữ liệu}

\subsection{InternetManager}

\codeword{InternetManager} cung cấp ba chế độ:
\codeword{INTERNET_MODE_NONE}, \codeword{INTERNET_MODE_WIFI} và
\codeword{INTERNET_MODE_MESH}. Khi đổi chế độ, module dừng các task mạng cũ,
đóng hoặc tái tạo netif, reset trạng thái mesh rồi mới khởi động stack mới.
Cờ \codeword{transitioning} ngăn hai thao tác chuyển chế độ chạy đồng thời.

\subsection{Hai vai trò mesh}

\begin{description}
  \item[Node:] tham gia Mesh-Lite, đọc dữ liệu cảm biến, tạo JSON và gửi UDP
  tới IP/port của root.
  \item[Root:] bind socket UDP, nhận frame từ node và đẩy vào
  \codeword{gateway_rx_queue}. Task UART lấy queue này và chuyển dữ liệu tới
  gateway ngoài.
\end{description}

Queue nhận tại root có độ sâu 64 frame; kích thước frame được giới hạn bởi
\codeword{MESH_ROOT_UDP_FRAME_SIZE}. Mutex bảo vệ việc công bố và vô hiệu hóa
socket khi đổi role hoặc dọn stack mạng.

\subsection{Nguồn dữ liệu và luồng telemetry}

Telemetry không chỉ chứa giá trị cảm biến. Gói hoàn chỉnh được tổng hợp từ
nhiều nguồn trong firmware: driver cảm biến cung cấp mảng \codeword{p}; bộ đếm
MeshManager cung cấp \codeword{seq}; Mesh-Lite cung cấp level \codeword{n};
MAC Wi-Fi STA tạo trường \codeword{M}; TimeManager cung cấp \codeword{t};
phiên bản firmware tạo \codeword{ver}; ErrorCodes cung cấp mảng
\codeword{err}. Sơ đồ sau chỉ rõ nguồn sinh từng nhóm trường và đường đi của
gói dữ liệu.

\begin{figure}[ht]
  \centering
  \begin{tikzpicture}[
    source/.style={draw=green!55!black,rounded corners=3pt,fill=green!8,
      minimum width=3.2cm,minimum height=0.9cm,align=center,font=\scriptsize},
    process/.style={draw=linkblue,rounded corners=3pt,fill=blue!8,
      minimum width=3.2cm,minimum height=1.05cm,align=center,font=\scriptsize},
    output/.style={draw=violet!70!black,rounded corners=3pt,fill=violet!8,
      minimum width=3.1cm,minimum height=1cm,align=center,font=\scriptsize},
    arrow/.style={-{Latex[length=2.5mm]},thick,draw=linkblue},
    label/.style={font=\tiny,fill=white,inner sep=1.5pt}
  ]
    \node[source] (sensor) {Sensor Registry + driver\\
      cổng, loại, các giá trị $\rightarrow$ \codeword{p}};
    \node[source,below=4mm of sensor] (identity) {Wi-Fi STA + Mesh-Lite\\
      MAC $\rightarrow$ \codeword{M}, level $\rightarrow$ \codeword{n}};
    \node[source,below=4mm of identity] (clock) {TimeManager + firmware\\
      thời gian $\rightarrow$ \codeword{t}, version $\rightarrow$ \codeword{ver}};
    \node[source,below=4mm of clock] (runtime) {MeshManager + ErrorCodes\\
      bộ đếm $\rightarrow$ \codeword{seq}, lỗi $\rightarrow$ \codeword{err}};

    \node[process,right=14mm of identity] (builder)
      {\textbf{MeshManager}\\đọc các nguồn và\\đóng gói JSON schema 1};
    \node[output,right=14mm of builder] (packet)
      {\textbf{Gói telemetry}\\
       \codeword{v,seq,n,M,t,ver,err,p}};

    \node[process,below=10mm of runtime] (root)
      {\textbf{Mesh root}\\UDP RX $\rightarrow$ queue\\tối đa 64 frame};
    \node[process,right=14mm of root] (normalize)
      {\textbf{ProcessingDataMesh}\\kiểm tra \codeword{M}\\tính packet loss};
    \node[output,right=14mm of normalize] (gateway)
      {UART gateway\\JSON + \codeword{\string\n}\\
       Backend / Dashboard};

    \draw[arrow] (sensor.east) -- (builder.west);
    \draw[arrow] (identity.east) -- (builder.west);
    \draw[arrow] (clock.east) -- (builder.west);
    \draw[arrow] (runtime.east) -- (builder.west);
    \draw[arrow] (builder) -- node[label,above]{tạo frame} (packet);
    \draw[arrow] (packet.south) |- node[label,near start,right]{UDP/Mesh}
      (root.east);
    \draw[arrow] (root) -- node[label,above]{frame nhận} (normalize);
    \draw[arrow] (normalize) -- node[label,above]{thay seq bằng packetloss}
      (gateway);
  \end{tikzpicture}
  \caption{Nguồn hình thành trường dữ liệu và tuyến truyền telemetry}
\end{figure}

Luồng xử lý được chia thành năm bước:

\begin{enumerate}
  \item SensorManager gọi callback \codeword{read()} của cảm biến đã gán cho
  từng cổng. Chỉ kết quả đọc thành công mới được thêm vào mảng
  \codeword{p}.
  \item MeshManager lấy MAC Wi-Fi STA làm định danh ổn định của node, lấy level
  hiện tại từ Mesh-Lite, timestamp từ TimeManager, phiên bản và danh sách lỗi.
  \item Các giá trị được ghép thành một JSON schema 1. Node gửi frame này bằng
  UDP tới địa chỉ root sau một khoảng jitter nhằm giảm xung đột khi nhiều node
  phát gần đồng thời.
  \item Root nhận UDP, giữ nguyên payload và đưa vào
  \codeword{gateway_rx_queue}. Nếu queue đầy, frame bị loại và bộ đếm
  \codeword{dropped_frames} tăng.
  \item ProcessingDataMesh dùng cặp \codeword{M}/\codeword{seq} để tính tỷ lệ
  mất gói của từng node trong cửa sổ 10 giây. Sau đó UART gửi một JSON trên mỗi
  dòng tới gateway, từ đó dữ liệu được chuyển tới backend và dashboard.
\end{enumerate}

\paragraph{Các nguồn dữ liệu phụ.}
BatteryManager, SystemPerformance, StorageManager và trạng thái mạng được lưu
trong DataManager để phục vụ màn hình hoặc chức năng chẩn đoán. Chúng không tự
động xuất hiện trong mảng \codeword{p}; muốn truyền thêm các chỉ số này phải
bổ sung trường vào schema, tăng phiên bản \codeword{v} nếu thay đổi không tương
thích và cập nhật parser phía gateway.

Lệnh UART \codeword{Start Root} yêu cầu thiết bị chuyển sang root nếu mesh đã
sẵn sàng. Gateway gửi heartbeat \codeword{Connected}; nếu quá một giây không
nhận heartbeat, firmware chuyển về node. Khi root không có frame trong queue,
firmware gửi \codeword{No node} định kỳ.

\section{Định dạng telemetry}

\subsection{Gói node hiện tại}

MeshManager tạo JSON theo thứ tự khóa cố định:
\codeword{v}, \codeword{seq}, \codeword{n}, \codeword{M}, \codeword{t},
\codeword{ver}, \codeword{err}, \codeword{p}.

\begin{lstlisting}[style=textstyle,caption={Ví dụ gói telemetry schema 1}]
{
  "v": 1,
  "seq": 281,
  "n": 2,
  "M": "B4:BF:E9:34:0C:CC",
  "t": "2012-01-01T00:06:23",
  "ver": "0.0.1",
  "err": ["0x51"],
  "p": [[1, 6, 27.142, 43.030]]
}
\end{lstlisting}

\begin{longtable}{p{1.8cm}p{3.2cm}p{9.5cm}}
  \caption{Ý nghĩa các trường telemetry}\\
  \toprule
  \textbf{Khóa} & \textbf{Kiểu} & \textbf{Ý nghĩa} \\
  \midrule
  \endfirsthead
  \toprule
  \textbf{Khóa} & \textbf{Kiểu} & \textbf{Ý nghĩa} \\
  \midrule
  \endhead
  \codeword{v} & số nguyên & Phiên bản schema; phải tăng khi thay đổi không
  tương thích. \\
  \codeword{seq} & số nguyên & Bộ đếm gói, dùng phát hiện mất gói. \\
  \codeword{n} & số nguyên & Level của node trong mạng mesh. \\
  \codeword{M} & chuỗi & MAC Wi-Fi STA, định danh nguồn phát. \\
  \codeword{t} & chuỗi & Thời gian ISO dạng YYYY-MM-DDTHH:MM:SS. \\
  \codeword{ver} & chuỗi & Phiên bản firmware major.minor.patch. \\
  \codeword{err} & mảng chuỗi & Các mã lỗi runtime dạng hexadecimal. \\
  \codeword{p} & mảng mảng & Mỗi hàng gồm cổng, loại cảm biến và các giá trị. \\
  \bottomrule
\end{longtable}

Một phần tử \codeword{p} có dạng
\codeword{[port, sensor_type, value_0, value_1, ...]}. Chỉ cổng đã gán cảm
biến và đọc thành công mới được đưa vào gói.

\subsection{Xử lý telemetry tại gateway}

Schema đã được thống nhất dùng khóa \codeword{M}. Hàm
\codeword{ProcessingDataMesh_ProcessFrame()} yêu cầu trường MAC này, dùng MAC
làm khóa riêng cho bảng trạng thái packet loss và giữ nguyên \codeword{M}
trong frame đầu ra. Gói không có \codeword{M} được xem là không hợp lệ thay vì
gộp nhầm vào một định danh mặc định. Bộ đệm định danh node có 18 byte, đủ chứa
17 ký tự MAC và ký tự kết thúc chuỗi.

\section{UI và tương tác người dùng}

UI gồm ButtonManager, ScreenManager và MenuSystem. Hai task độc lập thực hiện
đọc thao tác điều hướng và render màn hình. Trạng thái lựa chọn cảm biến của
ba cổng được lưu trong DataManager, vì vậy menu cấu hình và luồng thu thập dữ
liệu cùng sử dụng một nguồn trạng thái.

Khi mở rộng UI cần tránh thực hiện I/O chậm trong callback render. Thao tác
mạng, SD card hoặc cảm biến nên được chuyển cho task dịch vụ; UI chỉ đọc bản
chụp trạng thái và phát lệnh.

\section{Các dịch vụ lõi}

\begin{longtable}{p{4.2cm}p{10cm}}
  \caption{Vai trò các dịch vụ lõi}\\
  \toprule
  \textbf{Module} & \textbf{Trách nhiệm} \\
  \midrule
  \endfirsthead
  \toprule
  \textbf{Module} & \textbf{Trách nhiệm} \\
  \midrule
  \endhead
  BatteryManager & Đo điện áp, ước lượng mức pin và cập nhật nguồn cấp. \\
  TimeManager & Làm việc với RTC DS3231 và tạo timestamp cho telemetry. \\
  StorageManager & Ghi dữ liệu định kỳ; trên ESP32 được khởi tạo từ
  \codeword{app_main()}. \\
  SystemPerformance & Thu thập heap, tải CPU và các chỉ số runtime. \\
  ProcessingDataMesh & Parse JSON tại gateway, tính tỷ lệ mất gói theo cửa sổ
  10 giây và tạo frame đầu ra. \\
  UartToGateWay & Nhận lệnh từ gateway, quản lý heartbeat và chuyển frame của
  root ra UART. \\
  ErrorCodes & Chuẩn hóa mã lỗi và lưu lỗi vào mảng của DataManager. \\
  PinManager & Tập trung ánh xạ chân và cấu hình ngoại vi. \\
  \bottomrule
\end{longtable}

\section{Task FreeRTOS và tài nguyên}

\begin{table}[ht]
  \centering
  \caption{Một số task được tạo trực tiếp trong mã nguồn}
  \begin{tabularx}{\textwidth}{l c c X}
    \toprule
    Task & Stack & Priority & Chức năng \\
    \midrule
    \codeword{uart_gateway_rx} & 4096 & 6 & Nhận heartbeat và lệnh UART. \\
    \codeword{uart_gateway_tx} & 4096 & 7 & Chuyển queue mesh sang UART. \\
    \codeword{mesh_link} & 3072 & 5 & Theo dõi kết nối và socket mesh. \\
    \codeword{mesh_data} & 4096 & 5 & Gửi hoặc nhận telemetry UDP. \\
    Menu navigation & 4096 & 5 & Xử lý điều hướng menu. \\
    Menu render & 4096 & 5 & Cập nhật OLED. \\
    Battery task & 4096 & 5 & Đọc và cập nhật thông tin pin. \\
    Storage task & 4096 & 4 & Ghi dữ liệu định kỳ. \\
    System performance & 2048 & 5 & Thu thập chỉ số hệ thống. \\
    \bottomrule
  \end{tabularx}
\end{table}

Các con số stack là số đơn vị mà API ESP-IDF nhận cho task trong bản mã hiện
tại. Cần kiểm tra \codeword{uxTaskGetStackHighWaterMark()} trên thiết bị thật
trước khi giảm stack. Task UART TX có priority cao hơn để tránh queue root bị
đầy khi lưu lượng mesh tăng.

\section{Bộ nhớ flash và lưu trữ}

\begin{table}[ht]
  \centering
  \caption{Bảng phân vùng trong partitions.csv}
  \begin{tabular}{l l l l}
    \toprule
    Tên & Loại & Offset & Kích thước \\
    \midrule
    nvs & data/nvs & 0x9000 & 0x5000 \\
    otadata & data/ota & 0xe000 & 0x2000 \\
    ota\_0 & app/ota\_0 & 0x10000 & 0x1C0000 \\
    ota\_1 & app/ota\_1 & tự động & 0x1C0000 \\
    spiffs & data/spiffs & tự động & 0x20000 \\
    \bottomrule
  \end{tabular}
\end{table}

Hai phân vùng OTA cho phép cập nhật firmware theo mô hình A/B. SPIFFS có dung
lượng 128 KiB và chứa tài nguyên cấu hình web. Khi mount SPIFFS thất bại,
\codeword{init_spiffs()} hiện không tự format, giúp tránh mất dữ liệu ngoài ý
muốn nhưng yêu cầu người vận hành xử lý lỗi phân vùng.

\section{Build, flash và theo dõi}

\subsection{Yêu cầu}

\begin{itemize}
  \item ESP-IDF tương thích với các dependency của dự án;
  \item Python và toolchain Xtensa/RISC-V do ESP-IDF cung cấp;
  \item board ESP32 phù hợp với cấu hình target;
  \item cáp USB/UART và đúng cổng COM.
\end{itemize}

\subsection{Các lệnh cơ bản}

\begin{lstlisting}[style=textstyle,caption={Build và nạp firmware}]
cd MSMS_Project
idf.py set-target esp32
idf.py menuconfig
idf.py build
idf.py -p COMx flash monitor
\end{lstlisting}

Chỉ chạy \codeword{set-target} khi cần đổi chip vì thao tác này có thể tạo lại
cấu hình build. Các giá trị cần kiểm tra trong menuconfig gồm kênh mesh, IP
root, UDP port, baud UART gateway, Wi-Fi mesh và bảng phân vùng.

\subsection{Log cần theo dõi}

\begin{itemize}
  \item lỗi mount SPIFFS hoặc khởi tạo NVS;
  \item trạng thái layer mesh, parent RSSI và free heap;
  \item số frame UDP nhận, frame vào queue và frame bị drop;
  \item thống kê UART gồm frame, byte và mức sử dụng queue;
  \item lỗi đọc cảm biến và mảng \codeword{err} trong telemetry;
  \item reset reason, stack watermark và heap nhỏ nhất.
\end{itemize}

\section{Chiến lược kiểm thử}

\begin{longtable}{p{3.8cm}p{6.2cm}p{4.2cm}}
  \caption{Danh sách kiểm thử đề xuất}\\
  \toprule
  \textbf{Nhóm} & \textbf{Kịch bản} & \textbf{Kết quả mong đợi} \\
  \midrule
  \endfirsthead
  \toprule
  \textbf{Nhóm} & \textbf{Kịch bản} & \textbf{Kết quả mong đợi} \\
  \midrule
  \endhead
  Khởi động & Flash sạch, reset nóng, NVS lỗi phiên bản & Firmware khởi động,
  lỗi được ghi nhận, không treo. \\
  Cảm biến & Gắn từng loại, rút nóng, dữ liệu ngoài dải & Không crash; cổng lỗi
  không xuất dữ liệu sai. \\
  Mesh node & Mất parent rồi kết nối lại & Socket được phục hồi, sequence tiếp
  tục hợp lệ. \\
  Mesh root & Nhiều node gửi đồng thời & Queue không tràn trong tải thiết kế;
  thống kê drop chính xác. \\
  UART & Mất heartbeat \codeword{Connected} & Thiết bị chuyển từ root về node. \\
  Giao thức & Gói có \codeword{M}, thiếu \codeword{M}, sequence nhảy &
  Định danh đúng; gói thiếu MAC bị từ chối; packet loss đúng. \\
  Bộ nhớ & Chạy liên tục 24 giờ & Không giảm heap bất thường, không watchdog. \\
  OTA & Nâng cấp và rollback & Boot đúng partition, bảo toàn cấu hình cần thiết. \\
  \bottomrule
\end{longtable}

\section{Đánh giá và khuyến nghị}

\subsection{Điểm mạnh}

\begin{itemize}
  \item chia component rõ ràng và có registry để mở rộng cảm biến;
  \item tách vai trò node/root, có queue và thống kê drop;
  \item có lớp InternetManager để chuyển mode và dọn tài nguyên;
  \item có hai phân vùng OTA, quản lý lỗi, hiệu năng và pin;
  \item telemetry gọn, phù hợp thiết bị nhúng và băng thông thấp.
\end{itemize}

\subsection{Rủi ro kỹ thuật}

\begin{enumerate}
  \item DataManager là trạng thái toàn cục lớn; cần quy tắc rõ ràng về quyền
  ghi và đồng bộ giữa task.
  \item JSON được ghép bằng \codeword{snprintf}; phải kiểm thử sát giới hạn
  buffer khi tăng số trường hoặc số cảm biến.
  \item Một số comment nguồn có dấu hiệu sai mã hóa ký tự, gây khó bảo trì.
  \item Timestamp có thể bắt đầu ở năm 2012 nếu RTC chưa được đồng bộ; phía
  nhận cần đánh dấu dữ liệu thời gian chưa hợp lệ.
\end{enumerate}

\subsection{Thứ tự cải tiến đề xuất}

\begin{enumerate}
  \item thêm test tự động để khóa chặt schema telemetry dùng \codeword{M};
  \item định nghĩa rõ ownership và mutex của từng nhóm trường DataManager;
  \item bổ sung versioning và kiểm tra kích thước cho mọi frame;
  \item thêm test tải queue root với nhiều node và giám sát stack watermark;
  \item chuẩn hóa encoding toàn bộ file nguồn về UTF-8;
  \item đồng bộ tài liệu trong \projectpath{openwiki/} với chế độ mesh mặc định.
\end{enumerate}

\section{Kết luận}

\codeword{MSMS_Project} là firmware ESP-IDF theo kiến trúc component, hỗ trợ
nhiều cảm biến và hai vai trò trong mạng Mesh-Lite. Luồng chính đã tách tương
đối rõ giữa thu thập dữ liệu, giao diện, kết nối và chuyển tiếp gateway. Schema
định danh node đã được thống nhất dùng khóa \codeword{M}. Điểm cần ưu tiên
trước khi mở rộng hệ thống là bảo đảm truy cập DataManager an toàn giữa các
task và bổ sung kiểm thử hồi quy cho giao thức. Kiến trúc hiện tại có nền tảng
phù hợp để bổ sung cảm biến, OTA và các chức năng giám sát quy mô lớn hơn.

\appendix
\section{Tệp nguồn trọng tâm}

\begin{longtable}{p{14.2cm}}
  \toprule
  \textbf{Tệp nguồn và nội dung nên đọc} \\
  \midrule
  \endfirsthead
  \toprule
  \textbf{Tệp nguồn và nội dung nên đọc} \\
  \midrule
  \endhead

  \path{main/main.c}\\
  Trình tự khởi tạo hệ thống. \\[3pt]

  \path{component/core/DataManager/DataManager.h}\\
  Trạng thái trung tâm và ngữ cảnh mesh. \\[3pt]

  \path{component/sensors/SensorRegistry/SensorRegistry.c}\\
  Danh sách driver cảm biến. \\[3pt]

  \path{component/network/InternetManager/InternetManager.c}\\
  Chuyển đổi Wi-Fi/Mesh và dọn tài nguyên. \\[3pt]

  \path{component/network/MeshManager/MeshManager.c}\\
  Mesh, UDP, telemetry, queue và vai trò node/root. \\[3pt]

  \path{component/core/ProcessingDataMesh/ProcessingDataMesh.c}\\
  Parser telemetry và tính tỷ lệ packet loss theo MAC. \\[3pt]

  \path{component/core/UartToGateWay/UartToGateWay.c}\\
  Giao thức UART với gateway. \\[3pt]

  \path{partitions.csv}\\
  Bố trí flash và OTA. \\
  \bottomrule
\end{longtable}

\end{document}
